\par\indent\SourceLeaf{246}%
\textsc{Theorem.} \textit{If countless shortest lines of equal length are drawn on a curved surface from the same initial point, the line joining their endpoints is normal to each of them.}

We thought it worthwhile to deduce this theorem from the fundamental property of shortest lines. Its truth can also be seen without calculation by the following reasoning. Let \(AB,AB'\) be two shortest lines of the same length that enclose an infinitesimal angle at \(A\). Suppose one of the angles made by the element \(BB'\) with the lines \(BA,B'A\) differs from a right angle by a finite amount. By the law of continuity one angle will then be greater than a right angle and the other smaller. Suppose the angle at \(B\) is \(=90^0-\omega\). On the line \(BA\) take a point \(C\) such that \(BC=BB'\mathord{.}\mathrm{cosec}\,\omega\). Since the infinitesimal triangle \(BB'C\) can be treated as plane, we have \(CB'=BC\mathord{.}\cos\omega\), and therefore
\[
\begin{aligned}
AC+CB'&=AC+BC\mathord{.}\cos\omega=AB-BC\mathord{.}(1-\cos\omega)\\
&=AB'-BC(1-\cos\omega)
\end{aligned}
\]
that is, the path from \(A\) to \(B'\) through \(C\) is shorter than the shortest~line.~Q.~E.~A.
\Article{16}
We add another theorem to that of the preceding article and state it as follows. \textit{Suppose we take any line on a curved surface, with countless shortest lines of equal length leaving its individual points at right angles and toward the same side. The curve joining their other endpoints cuts each of them at right angles.} The preceding analysis proves this with one change. Here \(\varphi\) must denote the length of the \textit{given} curve measured from an arbitrary point or, if you prefer, a function of that length. All the reasoning still applies, with the modification that \(S=0\) for \(r=0\) is already implied by the hypothesis itself. This second theorem is more general than the preceding one, which can be regarded as contained in it when we take the given line to be an infinitesimal circle drawn about the center \(A\). Finally, we note that geometric considerations can replace the analysis here as well. They are evident enough that we need not dwell on them here.
\Article{17}
We return to the formula \(\BareRoot{(E\mathrm{d}p^2+2F\mathrm{d}p\mathord{.}\mathrm{d}q+G\mathrm{d}q^2)}\), which gives in general
